\section{The second-order term for two factors}
\label{sec:two-factor}

Consider a nonsplit extension $0\to F\to E\to G\to0$ with
$F,G$ nonisomorphic and stable of equal slope at $\omega_0$.
Write $a=\rk F$, $b=\rk G$ and $r=a+b$. In the stable side put
\[
 \delta_s=\kappa_{\omega_s}(G)-\kappa_{\omega_s}(F)>0,
 \qquad N_s=V_s\frac{ab}{r}.
\]
Theorem~\ref{thm:main} gives $\lambda_1\sim\delta_s$.
We now identify the next coefficient for the actual HE metrics.

\subsection{An exact identity at an HE metric}
Use the orthogonal smooth splitting determined by any HE metric on
$E$. Its Dolbeault operator has the form
\[
 \db_E=\begin{pmatrix}\db_F&\beta\\0&\db_G\end{pmatrix}.
\]
The diagonal metrics here are induced metrics; they need not be HE.
Let $P_F$ be the orthogonal projection onto $F$, and set
$q=P_F-(a/r)I$. Then $\|q\|^2=N_s$. The subbundle degree formula
gives
\begin{equation}\label{eq:exact-rayleigh}
 \|\beta\|^2=2\pi a\bigl(\mu(E)-\mu(F)\bigr)
             =N_s\delta_s,
 \qquad \frac{\langle q,\Delta q\rangle}{\|q\|^2}=\delta_s,
\end{equation}
because $\db_{\End}q=-\beta$. This Rayleigh upper bound is a
consequence of the classical degree identity.

The off-diagonal HE equation and the K\"ahler identity give
$\db_{\Hom}^*\beta=0$ for the induced connections. In a unitary
coframe write $\beta=\sum_\nu\beta_\nu\bar\theta^\nu$ and put
\begin{equation}\label{eq:C-beta}
 C_s(\beta)=\diag\left(\sum_\nu\beta_\nu\beta_\nu^*,
                       -\sum_\nu\beta_\nu^*\beta_\nu\right).
\end{equation}
Block multiplication in $\db_{\End}^*\db_{\End}q$ gives
\begin{equation}\label{eq:Delta-q}
 \Delta q=C_s(\beta).
\end{equation}
The term linear in $\beta$ is $-\db_{\Hom}^*\beta$ and vanishes;
the remaining quadratic term is \eqref{eq:C-beta}.

Let $Q$ project onto $q^\perp$ in complex trace-free endomorphisms,
and write $z=C_s(\beta)-\delta_sq$. Suppose the compression
$A=Q\Delta Q$ has spectrum at least $\gamma>\delta_s$.
For the unique small eigenvector, take its $q$ component to be one.
The complementary equation gives $v=-(A-\lambda)^{-1}z$.
Pairing the remaining equation with $q$ yields
\begin{equation}\label{eq:exact-schur-two}
 \lambda=\delta_s-\frac{\langle z,(A-\lambda)^{-1}z\rangle}{N_s}.
\end{equation}
In particular
\[
 \delta_s-\frac{\|z\|^2}{N_s(\gamma-\delta_s)}
       \le\lambda\le\delta_s.
\]

\subsection{The normalized extension limit}
Let $\beta_s$ be the harmonic representative for the split factor
HE metrics, and $B_s=\|\beta_s\|^2$. The scalar equation in
Section~\ref{sec:construction}, now paired with $q$, reads
\[
 -N_s\delta_s+B_st^2+O(t^4)=0.
\]
The coefficient $B_s$ is exact: the quadratic metric correction is
in the image of $L_s$ and pairs trivially with its parallel kernel
vector $q$. Parity makes the scalar equation even, and its derivative
in $w=t^2$ is $B_s+O(w)>0$. Consequently
\begin{equation}\label{eq:two-scaling}
 t_s^2=\frac{N_s}{B_s}\delta_s+O(\delta_s^2).
\end{equation}
These estimates are uniform in $s$, including flat approaches to
$\delta=0$.

We need an orthogonal identification that preserves the subbundle.
The involution $J=\diag(-I_F,I_G)$ sends $t$ to $-t$; uniqueness
gives $u(s,-t)=Ju(s,t)J$. Together with $u_t(s,0)=0$, this makes
the diagonal entries even of order $O(t^2)$ and the off-diagonal
entries odd of order $O(t^3)$. Write
\[
 e^u=\begin{pmatrix}A&B\\B^*&D\end{pmatrix},\qquad
 T=\begin{pmatrix}I&-A^{-1}B\\0&I\end{pmatrix}.
\]
Then
\[
 T^*e^uT=\diag(A,D-B^*A^{-1}B).
\]
Thus
$S=T\diag(A^{-1/2},(D-B^*A^{-1}B)^{-1/2})$ is an isometry
from the split metric to the corrected metric and preserves $F$.
Its diagonal correction is $O(t^2)$ and its upper entry is $O(t^3)$
in every $C^m$ norm. Including the derivative of $S$ in
$S^{-1}\db_{s,t}S$, the upper Dolbeault entry is therefore
\[
 \beta_{\mathrm{HE},s,t}=t\beta_s+O_{C^m}(t^3).
\]
The estimates at one higher spatial derivative justify this formula
in $C^m$. The diagonal connections converge smoothly to the split
wall connections. With \eqref{eq:two-scaling} we obtain
\begin{equation}\label{eq:normalized-beta}
 \frac{\beta_{\mathrm{HE},s}}{\sqrt{\delta_s}}
     \longrightarrow\beta_*:=\sqrt{\frac{N_0}{B_0}}\,\beta_0
 \quad\hbox{in every }C^m,\qquad \|\beta_*\|^2=N_0.
\end{equation}

This convergence takes place in the rescaled presentation of the
extension. On the original fixed $E$, the isomorphism to the scaled
extension multiplies $F$ by $t$; the pulled-back metric on $F$ is
therefore $t^2h_F$ to leading order. Smooth convergence to the split
metric is not asserted in that original presentation.

\subsection{The Green coefficient and eigensection}
Let $\Delta_0$ be the endomorphism Dolbeault Laplacian at the split
wall metric, and $G_0$ its inverse on the complement of $q_0$ in
trace-free endomorphisms. Define
\begin{equation}\label{eq:D-coefficient}
 z_0=C_0(\beta_*)-q_0,\qquad
 D=\frac{\langle z_0,G_0z_0\rangle}{N_0}.
\end{equation}
The vector $z_0$ is orthogonal to $q_0$, since
$\langle q_0,C_0(\beta_*)\rangle=\|\beta_*\|^2=N_0$.

\begin{theorem}\label{thm:second-order}
As $s\to0$ through $\delta_s>0$,
\begin{equation}\label{eq:second-order}
 \lambda_1(s)=\delta_s-D\delta_s^2+o(\delta_s^2).
\end{equation}
Here $D\ge0$, with equality precisely when $C_0(\beta_*)=q_0$
pointwise. After smooth unitary identifications with the wall spaces,
the normalized first eigensection, with positive $q_s$ component, is
\begin{equation}\label{eq:eigensection}
 \psi_s=\frac{q_s}{\sqrt{N_s}}
       -\delta_s\frac{G_0z_0}{\sqrt{N_0}}+o_{C^m}(\delta_s)
 \quad\hbox{for every }m.
\end{equation}
All other trace-free eigenvalues remain uniformly positive.
\end{theorem}

\begin{proof}
Use the orthogonal presentation above and an $L^2$ isometry $U_s$
incorporating the varying volume density. We use $q_s,Q_s,\Delta_s$
for the transported endomorphism objects. Nonlinear expressions must
be formed before this transport: in particular the transported
$C_s(\beta)$ means $U_sC_s(\beta)$, not $C_0$ applied to a
transported form. Density multiplication does not commute with the
quadratic operation $C$.

All differential coefficients converge smoothly to those of $\Delta_0$.
The latter has trace-free kernel $\CC q_0$. Its complementary gap
persists for $A_s=Q_s\Delta_sQ_s$. Indeed, a sequence of normalized
complementary sections with energy tending to zero would have bounded
$H^1$ norm by the first-order elliptic estimate, a strongly convergent
$L^2$ subsequence, and a nonzero limit in
$\ker\Delta_0\cap q_0^\perp$, a contradiction.
Min--max and \eqref{eq:exact-rayleigh} now give exactly one eigenvalue
$\lambda_s<\gamma/2$ for small $s$, and positivity follows from
simplicity of $E$.

Equations \eqref{eq:Delta-q} and \eqref{eq:normalized-beta} imply
\[
 z_s=Q_s\Delta_sq_s=U_sC_s(\beta_{\mathrm{HE},s})-\delta_sq_s,
 \qquad z_s/\delta_s\longrightarrow z_0
\]
in every Sobolev norm. After smooth fixed-complement identifications,
uniform elliptic estimates give
\[
 \|v\|_{H^{j+2}}\le C_j\bigl(\|(A_s-\lambda_s)v\|_{H^j}
                                      +\|v\|_{L^2}\bigr).
\]
The complementary gap absorbs the last term. The projections have
finite rank with smooth kernels and are uniformly bounded on every
Sobolev space. Hence
\[
 \|(A_s-\lambda_s)^{-1}\|_{H^j\to H^{j+2}}\le C_j.
\]
The resolvent identity and coefficient convergence from $H^{j+2}$
to $H^j$ show that these inverses converge to $G_0$ in the stated
operator norm. Substituting into \eqref{eq:exact-schur-two} proves
\eqref{eq:second-order}.

Before normalization the small eigensection is $q_s+v_s$, where
\[
 v_s=-(A_s-\lambda_s)^{-1}z_s
     =-\delta_sG_0z_0+o_{H^{j+2}}(\delta_s).
\]
Its squared norm is $N_s+\|v_s\|^2$. Normalization therefore changes
the $q_s$ coefficient only by $O(\delta_s^2)$. Sobolev embedding,
with arbitrary $j$, gives \eqref{eq:eigensection}; the phase is fixed
by $\langle q_s,\psi_s\rangle_s>0$. Positivity of $G_0$ proves
the assertion about $D$ and its zero set.
\end{proof}

The coefficient is unchanged by nonzero scalar rescaling of the
extension representative, after the induced unitary identifications.
It depends on the wall geometry and the projective extension class.
Ranks and Chern classes alone do not determine it. To compare with
\cite{Flow}, write $c=B_0/N_0$ and
$z=C_0(\beta_0)-cq_0$. Its fixed harmonic-family coefficient is
$d=\langle z,G_0z\rangle/N_0$. Since $z_0=z/c$, we have exactly
$D=d/c^2$. Theorem~\ref{thm:second-order} identifies this coefficient
for the actual HE metrics and the exact slope-gap parameter.
